\section{Related Work}
\label{sec:related}

We organize related work around three themes, concluding each with the limitation our work addresses.

\subsection{Perplexity-Based Quality Filtering}

Perplexity computed by a reference language model (typically GPT-2 or a domain-matched model) is a widely used proxy for training data quality. ProX~\cite{Zhou2024ProX} demonstrates that per-example programmatic text refinement, which includes a perplexity-informed quality signal, improves downstream benchmarks by approximately 2\% across C4, DCLM, and FineWeb corpora. REWIRE~\cite{Nguyen2025REWIRE} shows that transforming discarded low-quality documents (rather than simply filtering them) yields 1.0--2.5 percentage point improvements across 22 tasks at 1--7B scale.

Critically, DataMan~\cite{Peng2025DataMan} identifies a misalignment between perplexity-based quality and downstream in-context learning (ICL) performance: the token mixture that minimizes pre-training perplexity is not the same mixture that maximizes ICL scores. This PPL/ICL misalignment motivates careful threshold selection. However, DataMan evaluates this misalignment at a single model scale. Our work extends this line of inquiry by testing whether the optimal threshold in the PPL/ICL trade-off shifts with model scale --- finding that it does, and dramatically so.

\subsection{Deduplication Strategies}

Near-duplicate removal is a standard component of pre-training data pipelines. SoftDedup~\cite{He2024SoftDedup} shows that $n$-gram commonness-based soft weighting of near-duplicates improves few-shot performance by 1.77\% over hard deduplication on a fixed-scale model. FineWeb2~\cite{Penedo2025FineWeb2} provides principled multi-language curation ablations including deduplication tuning, reporting improved benchmark performance across languages.

Both works evaluate deduplication effects at a single model scale. Our experiment includes two deduplication levels (MinHash Jaccard $J{=}0.7$ strict, $J{=}0.9$ loose) as a secondary axis in the factorial design. While our primary gate focused on the PPL $\times$ scale interaction, the experimental data exists for a post-hoc dedup $\times$ scale analysis --- a direction we identify as immediate future work.

\subsection{Scalable Curation Ablation Methodology}

The computational expense of from-scratch pre-training ablations motivates work on efficient approximation. \citet{Na2024Scalable} propose training small proxy models and merging checkpoints to approximate ablation results, reporting Spearman $r{=}0.81$ correlation between proxy-model performance and full-scale benchmark outcomes. Chinchilla scaling laws~\cite{Hoffmann2022Training} provide theoretical grounding for relating compute-optimal training configurations across scales.

Our h-e1 proof-of-concept experiment (7M/16M proxy models, 200 training steps) produced $p{=}1.0$ and $\eta^2{\approx}0$ with all models at the MMLU floor --- no statistical signal. This contradicts straightforward application of \citeauthor{Na2024Scalable}'s proxy-model approach for Scale $\times$ Curation interaction studies: in our experiments, below a minimum capacity threshold (provisionally, approximately 14M parameters with a 2.2$\times$ scale ratio at 200--500 training steps), the interaction cannot manifest because neither model is large enough to exhibit the capacity-quality trade-off behavior. We note that this threshold may depend on scale ratio and training duration rather than absolute parameter count alone.

\subsection{Our Position}

To our knowledge, we contribute the first controlled factorial experiment that directly tests the Scale $\times$ Curation interaction in LLM pre-training. Prior multi-scale curation work, including FineWeb2~\cite{Penedo2025FineWeb2} and DataComp-LM~\cite{Gururangan2024DataComp}, conducts ablations at multiple scales for the purpose of tuning a single unified curation recipe --- these are not designed as factorial tests of the scale-curation \emph{interaction} (i.e., whether the optimal recipe differs by scale). Unlike prior work, which either (a) evaluates curation at a single scale, (b) assumes scale-independence, or (c) uses proxy models too small to exhibit scale-dependent effects, we run a 24-run factorial experiment at 14M/31M parameter scales with real FineWeb data, real GPT-2 PPL scoring, and standard evaluation. The resulting $\tau^*(14\text{M}){=}20$ vs.\ $\tau^*(31\text{M}){=}50$ finding provides empirical evidence that optimal curation thresholds are model-scale-dependent at PoC scale.
